\input{1-preamble}

\title{When Allies Have Your Back:\\
Allied Support and Public Responses to Coercive Threats}

\author{Hannah Jakob Barrett}

\begin{document}
\maketitle

\textit{Pre-Analysis Plan --- Workshop Draft. Please do not cite or circulate. Comments welcome, especially on the vignette design and case selection.}

\newpage

%-------------------------------------------------------------------
\section{Introduction}
%-------------------------------------------------------------------

The international system is replete with coercive threats: demands backed by the implicit or explicit threat of military force. How the public in a target state responds to such threats has significant consequences for the decisions that democratic leaders ultimately take. If citizens categorically oppose compliance, leaders face domestic pressure to stand firm even when the strategic calculus might favour concession. If, conversely, the public views compliance as acceptable, leaders retain room to defuse crises through diplomacy.

The experimental literature on public responses to coercive threats has made important advances in documenting the psychological mechanisms that drive defiance. Studies conducted predominantly in the United States find consistent evidence of coercive failure: rather than prompting compliance, threats tend to increase public support for retaliation through mechanisms of reactance and reputational concern.\autocite{dafoe_coercion_2021, dafoe_provocation_2022, powers_psychology_2023} Citizens become angry and resistant when faced with an adversarial ultimatum; they interpret compliance as a humiliating loss of national face. Yet this literature suffers from a systematic limitation. The overwhelming majority of experiments test responses among American citizens---the citizens of the world's foremost military superpower. Respondents confident in their country's overwhelming military superiority are predisposed toward defiance regardless of the threat's severity. The mechanisms of coercive failure documented in US samples may not generalise to the population of states that more commonly face such threats: small and medium-sized states operating under the shadow of more powerful adversaries.

This pre-analysis plan addresses this gap by studying public responses to a coercive threat among Danish citizens---a small NATO member with direct exposure to Russian military pressure in the Baltic region. By situating the adversary as a clearly more powerful state (Russia), the design tests whether the coercive failure mechanisms documented in US samples hold when the costs of defiance are objectively higher.

More critically, this study introduces a theoretically and empirically underexplored variable: the role of allied support in shaping public responses to coercive threats. Whether and how allied backing shapes the perceived costs of compliance versus defiance is fundamental to understanding how alliances function in crisis situations. Yet direct experimental tests of this mechanism remain rare, and no study has examined whether the \textit{identity} of the committing ally---European partners versus the United States---conditions the emboldening effect. In the current geopolitical context, where the credibility of US security commitments to European allies has come under sustained scrutiny, this question is both theoretically important and practically urgent.

I propose a five-condition survey experiment varying the nature of allied support in response to a Russian maritime ultimatum directed at Denmark. The design distinguishes between no allied commitment, European-only commitment, US-only commitment, full NATO commitment, and a restraining condition in which allies urge compliance. This allows me to estimate the main effect of allied support on defiance, test whether the quality and identity of the committing actor conditions that effect, and assess whether allied calls for restraint reduce public support for defiance.

The study makes three contributions. First, it provides direct experimental evidence on the emboldening versus restraining effects of alliances on public opinion---a mechanism widely theorised but rarely tested at the level of mass publics. Second, by disaggregating allied support by actor identity, it generates evidence on the comparative credibility of US versus European security commitments to European publics---a question of immediate policy relevance given recent transatlantic tensions. Third, it extends the coercion literature beyond US samples to a small-state context in which the risk calculus fundamentally differs.


%-------------------------------------------------------------------
\section{Theory}
%-------------------------------------------------------------------

\subsection{Coercive Threats and the Costs of Compliance}

Compellent threats---demands that a target state change its behaviour under threat of force---impose a fundamental asymmetry on the target: compliance requires a publicly visible capitulation, while defiance preserves the status quo at the risk of military escalation.\autocite{schelling_arms_1966} This asymmetry helps explain why compellence is empirically harder to achieve than deterrence.\autocite{schelling_arms_1966, sechser_goliaths_2010} Compliance with an ultimatum is not merely a strategic concession; it is a reputational event. A state that yields to coercive pressure signals to the adversary---and to third parties---that further threats will be effective. The reputational costs of compliance therefore extend beyond the immediate crisis: capitulation risks inviting future demands and undermining the state's broader deterrent posture.\autocite{fearon_domestic_1994, dafoe_coercion_2021} Resistance may therefore be rational not solely because of honour, but because compliance creates a precedent of vulnerability to coercion.

\textcite{sechser_goliaths_2010} shows that even powerful states frequently fail to extract compliance through coercive threats, precisely because targets resist accepting the reputational costs of visible capitulation. This dynamic operates at the level of public opinion as well. \textcite{dafoe_coercion_2021, dafoe_provocation_2022} demonstrate experimentally that coercive threats activate honour and reputational concerns among ordinary citizens, increasing public support for retaliation as a means of protecting national face. The public's resistance is not irrational; it reflects a coherent calculation that compliance with an ultimatum sets a precedent that makes future coercion more likely.

I argue that this reputation mechanism produces a \textit{default} disposition toward defiance in response to coercive threats. Citizens who learn that their country has been issued an ultimatum and threatened with military force will---absent other information---tend to favour standing firm. Compliance requires accepting a visible loss; defiance preserves national reputation at the cost of escalation risk. The key question is what information shifts the perceived costs of defiance sufficiently to alter this default preference.

\subsection{Reactance and Emotional Responses to Coercive Threats}

A second pathway to defiance operates through emotional, rather than strictly strategic, mechanisms. Building on reactance theory from psychology, \textcite{powers_psychology_2023} show that coercive threats increase public support for retaliation by framing the adversary's demand as an infringement on the target state's freedom of action. When citizens perceive their country as being told what to do under threat of violence, they experience autonomy threat---a psychological state that produces resistance and anger independent of any strategic calculation about costs and benefits.

This emotional mechanism complements, rather than competes with, the reputational mechanism described above. Reputation theory predicts defiance because compliance is costly;\autocite{dafoe_coercion_2021} reactance theory predicts defiance because coercion is resented.\autocite{powers_psychology_2023} Anger is associated with greater support for retaliation against external threats,\autocite{shandler_cyber_2021, fisk_emotions_2019} while fear and anxiety are associated with more cautious, de-escalatory preferences.\autocite{huddy_threat_2005} Coercive threats, precisely because they are perceived as aggressive and constraining, are more likely to activate anger than fear---reinforcing the default disposition toward defiance through an emotional channel.

The key implication is that citizens who favour defiance may do so for both rational (contagion risk, reputation) and emotional (reactance, anger) reasons. Both mechanisms can be probed using dedicated survey batteries, allowing this study to shed light on the relative weight of strategic versus emotional drivers of public opinion---even if the primary causal question concerns the role of allied support.

\subsection{Allied Support as a Modulator of Compliance Costs}

Military alliances affect the costs of defiance in two theoretically distinct ways: through an \textit{emboldening} effect and a \textit{restraining} effect.\autocite{fang_concede_2014} The emboldening effect holds that allied military backing reduces the perceived risks of defiance. A target state with allies committed to its defence can stand firm against a coercive ultimatum at lower cost, since the risks of military escalation are distributed across the alliance. This risk-sharing function of alliances is central to their deterrent value.\autocite{leeds_alliances_2003, benson_inducing_2014} I argue that this effect operates at the level of public opinion: citizens who believe their country's allies have committed to its defence should update their perceptions of the risks of defiance downward, increasing support for standing firm.

The restraining effect operates in the opposite direction. When allies recommend that the target state comply with an adversary's demand---to avoid escalation, preserve alliance unity, or protect broader strategic interests---deviating from that advice imposes reputational costs within the alliance. Defying allied counsel signals unreliability as a partner and risks the withdrawal of future support.\autocite{fang_concede_2014} I argue that public audiences are sensitive to these intra-alliance reputational costs: citizens whose allies recommend compliance may reduce their support for defiance out of concern for the country's standing within the alliance. Reckless retaliation that alienates committed allies may thus be perceived as strategically counterproductive even among citizens predisposed toward standing firm.

These two effects generate competing predictions depending on the nature of the allied signal. When allies commit to supporting defiance, they reduce the costs of standing firm and should increase public support for retaliation. When allies counsel compliance, they raise the reputational costs of defiance and should reduce public support for standing firm.

\paragraph{The quality and identity of allied commitment.}
Not all allied commitments are equally credible, and not all allies are equally capable. I argue that both the identity of the committing ally and the comprehensiveness of the commitment shape how much the public updates its assessment of the costs of defiance. A commitment from the United States---the alliance's most militarily capable member---provides a more credible guarantee than a commitment from smaller European partners alone, because the US possesses the capacity to decisively affect the outcome of a military confrontation. Full NATO commitment, backed by both the US and European members, should therefore produce the strongest emboldening effect.

However, the credibility of US security commitments to European allies has come under sustained pressure since the Trump administration's repeated questioning of NATO's collective defence obligation. If Danish citizens have internalised doubts about the reliability of US guarantees, a US-only commitment may generate a smaller emboldening effect than the quality of US military capability would otherwise predict. This is an empirically open question with direct policy implications: if the effect of a US commitment on public support for defiance is indistinguishable from zero---or substantially smaller than the effect of European-only commitment---this constitutes direct evidence of US credibility decay among European publics.

\paragraph{Allied support and secondary attitudes.}
Beyond its effects on primary crisis behaviour, allied support may affect citizens' broader assessments of the alliance. Citizens who observe their allies committing to their defence may update positively on confidence in NATO and trust in the alliance as an institution. Conversely, citizens in the restraining condition---whose allies urge compliance rather than offering support---may update negatively on both dimensions if they interpret the allied recommendation as abandonment in a moment of need.


%-------------------------------------------------------------------
\section{Research Hypotheses}
%-------------------------------------------------------------------

Table \ref{tab:hypotheses} presents the pre-registered hypotheses. Confirmatory hypotheses are tested at $\alpha = 0.05$ using the Holm-Bonferroni step-down procedure to correct for multiple comparisons across the primary tests (see Section \ref{sec:model}). Exploratory hypotheses are not subject to this correction and are interpreted as hypothesis-generating.

\begin{table}[H]
\centering
\caption{Confirmatory and Exploratory Hypotheses}
\label{tab:hypotheses}
\footnotesize
\begin{tabular}{p{1.2cm} p{5.5cm} p{5.3cm} p{1.2cm}}
\toprule
\textbf{Label} & \textbf{Hypothesis} & \textbf{Proposed Test} & \textbf{Type} \\
\midrule
H1 & Allied support (T1, T2, T3) increases support for defiance relative to no commitment (C). & T1 vs.\ C; T2 vs.\ C; T3 vs.\ C (planned contrasts, Holm-Bonferroni corrected). & Confirmatory \\[8pt]
H2 & The emboldening effect increases with the comprehensiveness and capability of allied commitment: T3 $>$ T1 and T3 $>$ T2. & T3 vs.\ T1; T3 vs.\ T2 (planned contrasts). & Confirmatory \\[8pt]
H3 & Allied restraint (T4) decreases support for defiance relative to no commitment (C). & T4 vs.\ C (planned contrast). & Confirmatory \\[8pt]
\midrule
H4 & US credibility decay: the effect of US-only commitment (T2) is smaller than the effect of European-only commitment (T1). & T2 vs.\ T1. & Exploratory \\[8pt]
H5 & Hawkish respondents (high militant internationalism) show larger emboldening effects of allied support. & Treatment $\times$ Militant Internationalism interaction. & Exploratory \\[8pt]
H6 & Anger partially mediates the effect of allied support on defiance. & Causal mediation analysis with anger index as mediator. & Exploratory \\[8pt]
H7 & Allied support (T1--T3) increases confidence in NATO allies and trust in NATO; allied restraint (T4) decreases both. & OLS regressions with NATO confidence and NATO trust as dependent variables. & Exploratory \\[8pt]
H8 & Allied restraint (T4) activates reputational concerns about deviating from allied recommendations, reducing defiance through this channel. & Causal mediation with reputation index as mediator, T4 vs.\ C. & Exploratory \\
\bottomrule
\end{tabular}
\end{table}


%-------------------------------------------------------------------
\section{Research Design}
\label{sec:design}
%-------------------------------------------------------------------

\subsection{Sample and Recruitment}

Participants will be recruited from YouGov's Danish online panel. The target sample is $n \approx 3{,}200$ voting-eligible adults (age 18 and above), allocated equally across the five treatment conditions ($n \approx 640$ per condition). Respondents who fail attention checks will be excluded from the analytical sample prior to hypothesis testing; target cell sizes refer to the post-exclusion sample.

\subsection{Experimental Design}

The study employs a single-factor between-subjects design with five conditions. All respondents read an identical crisis scenario. The experimental manipulation is contained in the final paragraph of the vignette, which varies the nature---or absence---of an allied statement in response to the Russian ultimatum. Table \ref{tab:design} presents the five conditions.

\begin{table}[H]
\centering
\caption{Five-Condition Design}
\label{tab:design}
\footnotesize
\begin{tabular}{p{0.4cm} p{3.2cm} p{8.8cm}}
\toprule
 & \textbf{Condition} & \textbf{Allied Signal} \\
\midrule
C  & No commitment       & Denmark's NATO allies have not reached a formal consensus and have not issued a collective statement about the incident. \\[5pt]
T1 & European commitment & NATO's European member states have issued a joint statement formally committing to provide Denmark with military support should the conflict escalate. \\[5pt]
T2 & US commitment       & The United States has issued a formal statement committing to provide Denmark with military support should the conflict escalate. \\[5pt]
T3 & Full NATO commitment & NATO has issued a formal collective statement, endorsed by all member states including the United States, committing to provide Denmark with military support should the conflict escalate. \\[5pt]
T4 & Allied restraint    & Denmark's NATO allies have issued a joint statement urging Denmark to release the vessel and resolve the dispute through diplomatic channels, warning that military escalation would not serve the alliance's collective interests. \\
\bottomrule
\end{tabular}
\end{table}

\subsection{Vignette}

The crisis scenario is presented as a short breaking-news article, formatted to resemble a wire service report. The article is split across three consecutive survey pages to encourage attentive reading; respondents click through to advance. The treatment paragraph (the allied statement) appears on the third page.

\bigskip
\noindent\rule{\textwidth}{0.4pt}
\smallskip

\noindent {\small \textsc{International News} \hfill \textit{Breaking}}

\medskip

\noindent \textbf{Russia Issues Naval Ultimatum to Denmark}

\medskip

\noindent \textit{Page 1 --- What happened:} Danish naval forces have detained a Russian vessel in the Baltic Sea after it was found shadowing a Danish patrol ship with its transponder disabled. The vessel was boarded and its crew taken into custody pending investigation.

\medskip

\noindent \textit{Page 2 --- Russia's response:} Russian authorities have denied any wrongdoing. In a public statement, the Russian government issued the following demand: if Denmark does not immediately release the vessel and crew, Russia will not hesitate to take military action to secure their release. The statement was broadcast on international news.

\medskip

\noindent \textit{Page 3 --- What allies are saying:} \textit{[Allied treatment text --- see Table \ref{tab:design}]}

\smallskip
\noindent\rule{\textwidth}{0.4pt}
\bigskip

\subsection{Manipulation Checks}

Following the vignette and prior to the outcome questions, respondents complete four manipulation checks administered in randomised order:

\begin{enumerate}
    \item \textit{Ultimatum check:} In the scenario you just read, did Russia issue a direct demand to Denmark?
    \begin{enumerate}
        \item Yes
        \item No
        \item They requested something but did not demand it
        \item I don't know
    \end{enumerate}
    \item \textit{Power check:} Based on the scenario, which country has superior military capabilities?
    \begin{enumerate}
        \item Russia is stronger
        \item They are roughly equal
        \item Denmark is stronger
        \item I don't know
    \end{enumerate}
    \item \textit{Allied commitment check:} How confident are you that Denmark's allies would provide military support if this conflict were to escalate? (slider, 0\% -- 100\%)
    \item \textit{Allied signal check:} Based on the scenario, what position did Denmark's allies take?
    \begin{enumerate}
        \item They committed to supporting Denmark militarily
        \item They urged Denmark to comply with Russia's demand
        \item They took no collective position
        \item I don't know
    \end{enumerate}
\end{enumerate}

Respondents who fail manipulation checks will not be excluded from the primary analysis. A pre-specified robustness check will replicate the main results on the subset of respondents who pass all four checks.


%-------------------------------------------------------------------
\section{Measurement}
\label{sec:measurement}
%-------------------------------------------------------------------

\subsection{Primary Outcome: Comply or Defy}

The primary dependent variable confronts respondents with a direct binary choice:

\begin{quote}
\textit{What do you think Denmark should do in response to Russia's ultimatum?}

\quad 1 --- Release the vessel and crew immediately \hfill [Comply]

\quad 2 --- Continue holding the vessel \hfill [Defy]
\end{quote}

\noindent The item is coded 0 (Comply) and 1 (Defy). Following this choice, respondents indicate the strength of their preference:

\begin{quote}
\textit{On a scale from 0 to 10, how strongly do you feel about this choice?}
\end{quote}

\subsection{Escalation Scale}

All respondents are then asked:

\begin{quote}
\textit{If Denmark were to continue holding the vessel, to what extent would you support or oppose each of the following approaches? Please rate each option on a scale from ``Strongly support'' to ``Strongly oppose.''} (Presented in randomised order; 5-point scale)
\end{quote}

\begin{enumerate}
    \item Pursue de-escalation through diplomatic channels, without taking any military action
    \item Impose economic sanctions on Russia
    \item Deploy additional Danish naval forces to the area to confront Russia's vessels
    \item Conduct a military strike against Russian naval forces
\end{enumerate}

These four items constitute an ordered escalation scale spanning de-escalatory diplomacy through military force. They are analysed both as individual items and as a composite mean rescaled to 0--1. The scale is qualitatively distinct from the primary binary in that it measures \textit{how much} defiance the respondent favours and at what cost: the de-escalation item captures respondents who favour holding firm while minimising escalation risk; the military strike item captures those willing to pay the highest costs of resistance.

\subsection{Open-Ended Response}

Following the escalation scale, respondents are asked:

\begin{quote}
\textit{In a few sentences, please explain why you think Denmark should follow that approach.}
\end{quote}

Open-ended responses will be manually coded into thematic argument categories following the procedures described in \textcite{jakob_barrett_how_2026}.

\subsection{Mechanism Measures}

\paragraph{Emotions battery.}
\textit{Please indicate your feelings toward Russia when thinking about this situation. To what extent do you feel\ldots} (presented in randomised order; 0 = not at all, 10 = very much)

\begin{enumerate}
    \item Angry \hfill \textit{[anger]}
    \item Furious \hfill \textit{[anger]}
    \item Annoyed \hfill \textit{[anger]}
    \item Anxious \hfill \textit{[fear]}
    \item Worried \hfill \textit{[fear]}
    \item Sad \hfill \textit{[distractor]}
    \item Happy \hfill \textit{[distractor]}
\end{enumerate}

An anger index and a fear index are computed as the mean of the respective items.

\paragraph{Reputation battery.}
\textit{To what extent do you agree or disagree with the following statements?} (presented in randomised order; 5-point scale from \textit{Strongly disagree} to \textit{Strongly agree})

\begin{enumerate}
    \item This situation puts Denmark's international reputation at stake.
    \item Russia has humiliated Denmark.
    \item Releasing the vessel now would invite further Russian threats to Denmark's security in the long run.
    \item Releasing the vessel would reduce the risk of war with Russia. \hfill \textit{[reverse-coded]}
\end{enumerate}

\paragraph{Reactance battery.}
\textit{To what extent do you agree or disagree with the following statements?} (presented in randomised order; 5-point scale from \textit{Strongly disagree} to \textit{Strongly agree})

\begin{enumerate}
    \item Russia tried to make a decision for Denmark.
    \item Russia tried to pressure Denmark into acting against its own interests.
    \item Russia tried to take away Denmark's freedom to choose how to respond.
\end{enumerate}

\subsection{Secondary Outcomes}

\paragraph{National resolve.}
\textit{If this conflict were to escalate, how confident are you that Denmark would stand firm against Russia's ultimatum?} (0--10 scale)

\paragraph{NATO confidence.}
\textit{How confident are you that Denmark's NATO allies would come to its defence if Denmark faced a serious military threat?} (0--10 scale)

\paragraph{NATO trust.}
\textit{How much do you trust NATO as an institution to protect Denmark's security?} (0--10 scale)

\subsection{Pre-Treatment Measures}

\begin{enumerate}
    \item \textbf{Demographics:} age, gender, education level
    \item \textbf{Partisanship:} prospective vote intention (party list)
    \item \textbf{Political interest:} \textit{How interested are you in politics?} (1 = Not at all interested, 5 = Very interested)
    \item \textbf{Foreign policy postures} \autocite{rathbun_towards_2020}. Respondents rate their agreement with each item on a 5-point scale from \textit{Strongly disagree} to \textit{Strongly agree}:
    \begin{enumerate}
        \item \textit{Isolationism 1:} Denmark's interests are best protected by avoiding involvement with other nations.
        \item \textit{Isolationism 2:} Denmark needs to simply mind its own business when it comes to international affairs.
        \item \textit{Militant internationalism 1:} Denmark should take all steps including the use of force to prevent aggression by any expansionist power.
        \item \textit{Militant internationalism 2:} Denmark needs a strong military to be effective in international relations.
        \item \textit{Cooperative internationalism 1:} When deciding its foreign policies, Denmark should take into account the views of its major allies.
        \item \textit{Cooperative internationalism 2:} Denmark should be more committed to diplomacy and not so fast to use the military in international crises.
    \end{enumerate}
    \item \textbf{Pre-treatment attitude toward Russia:} \textit{On a scale from 0 to 10, how favourable is your overall opinion of Russia?}
    \item \textbf{Pre-treatment attitude toward NATO:} \textit{On a scale from 0 to 10, how favourable is your overall opinion of NATO?}
\end{enumerate}


%-------------------------------------------------------------------
\section{Statistical Model}
\label{sec:model}
%-------------------------------------------------------------------

\subsection{Primary Analysis}

The primary analysis estimates the average treatment effect of each condition relative to the no-commitment control (C). Let $Y_i \in \{0, 1\}$ denote the binary comply/defy outcome for respondent $i$, where $1 =$ Defy. The estimating equation is:

\begin{equation}
    Y_i = \alpha + \beta_1 \text{T1}_i + \beta_2 \text{T2}_i + \beta_3 \text{T3}_i + \beta_4 \text{T4}_i + \varepsilon_i
    \label{eq:main}
\end{equation}

where $\text{T1}_i$--$\text{T4}_i$ are binary indicators for treatment assignment and C is the omitted reference category. Equation \eqref{eq:main} is estimated via OLS (linear probability model) with HC2 heteroskedasticity-robust standard errors. Logistic regression is reported as a robustness check; average marginal effects are reported at the mean of all covariates.

A covariate-adjusted specification includes pre-treatment controls (age, gender, foreign policy postures, pre-treatment NATO attitude) to improve precision. In a randomised experiment, treatment is independent of pre-treatment covariates by design; covariate-adjusted estimates remain unbiased and are included solely for efficiency gains.

The same specification is estimated with the continuous escalation scale as the dependent variable using OLS with robust standard errors.

\paragraph{Confirmatory contrasts and multiple testing.}
The confirmatory hypotheses correspond to the following planned contrasts: (H1) $\hat{\beta}_1 > 0$, $\hat{\beta}_2 > 0$, $\hat{\beta}_3 > 0$; (H2) $\hat{\beta}_3 > \hat{\beta}_1$ and $\hat{\beta}_3 > \hat{\beta}_2$; (H3) $\hat{\beta}_4 < 0$. These contrasts are evaluated using the Holm-Bonferroni step-down procedure applied to the family of primary contrasts at $\alpha = 0.05$. Exploratory hypotheses (H4--H8) are not subject to multiplicity correction.

\subsection{Heterogeneous Effects}

Exploratory heterogeneous effects (H5) are estimated by interacting treatment indicators with the pre-treatment militant internationalism scale:

\begin{equation}
    Y_i = \alpha + \sum_{k=1}^{4} \beta_k \text{T}k_i + \gamma \cdot \text{MilInt}_i + \sum_{k=1}^{4} \delta_k \bigl(\text{T}k_i \times \text{MilInt}_i\bigr) + \varepsilon_i
\end{equation}

The $\delta_k$ coefficients test whether the treatment effect of condition $k$ is larger among hawkish (high militant internationalism) respondents. Analogous interaction specifications are estimated with cooperative internationalism and isolationism as moderators.

\subsection{Mediation Analysis}

The causal mediation analyses (H6, H8) follow \textcite{imai_general_2010} as implemented in the \texttt{mediation} package in R.\autocite{tingley_mediation_2014} For H6, the anger index is specified as the mediator in each treatment-versus-control comparison. For H8, the reputation index is specified as the mediator in the T4 versus C comparison. Average causal mediation effects (ACME) and average direct effects (ADE) are estimated with 1{,}000 bootstrapped replications. Results are treated as exploratory.


%-------------------------------------------------------------------
\section{Power Analysis}
%-------------------------------------------------------------------

\textit{To be completed following finalisation of the experimental design and pre-registration. Target: 80\% power to detect an effect of $d = 0.10$ on the binary primary outcome at $\alpha = 0.05$, two-sided, for the primary contrast T3 vs.\ C with $n \approx 640$ per condition. Power calculations will be conducted in R.}


%-------------------------------------------------------------------
\section{Case Selection}
%-------------------------------------------------------------------

Denmark is selected as the primary case for three reasons. First, as a small NATO member with direct exposure to Russian military pressure in the Baltic Sea, Denmark exemplifies the class of states for which the theory is most relevant: states that face a materially stronger adversary and whose security depends substantively on alliance credibility. The naval scenario is ecologically valid and does not require respondents to reason about an unfamiliar threat context. Second, Denmark's multiparty political system provides sufficient within-sample variation in foreign policy orientations to test heterogeneous treatment effects. Third, studying a non-US, non-great-power sample directly addresses the sample selection bias in the existing experimental coercion literature.

The survey instrument names Denmark as the target country and Russia as the adversary. This choice prioritises ecological validity and allows the results to speak directly to a realistic policy scenario, at the cost of some internal validity with respect to respondents' pre-existing beliefs about Russia's capabilities and intentions. Extensive pre-treatment measures and the randomisation of treatment assignment mitigate this concern.

A French comparative sample is considered as a secondary case, contingent on budget. France's Gaullist strategic tradition emphasises autonomy from NATO and, particularly, from the United States, suggesting that the emboldening effect of allied support---especially of US commitment---should be attenuated relative to Denmark. A Denmark--France comparison would thus provide a between-country test of strategic culture as a moderator of the allied support effect, and would offer a more demanding test of the US credibility decay hypothesis (H4).

\printbibliography

\end{document}
